% title: Divisors with Prime or Dividing Successor
% short-title: Successor Divisor Problem
% abstract: We formalize in Lean 4 the divisor condition asking for all positive integers \(n\) such that for every positive divisor \(d\mid n\), either \(d+1\mid n\) or \(d+1\) is prime. The proof is structural: after separating the odd part of \(n\), the hypothesis forces that odd part to be one less than a power of two, and a second divisor argument reduces the remaining cases to \(n\in\{1,2,4,12\}\).
% subjclass: 03B35
% keywords: divisors, number theory, olympiad problem, Lean 4
% author: Mario Hernández M.

\subsection{Divisors with Prime or Dividing Successor}

We define the property
\lean{Biblioteca.Demonstrations.SuccessorDivisorOrPrime} and prove the
classification theorem
\lean{Biblioteca.Demonstrations.successor_divisor_or_prime_classification}.

\begin{theorem}
The positive integers \(n\) such that for every positive divisor \(d\mid n\),
either \(d+1\mid n\) or \(d+1\) is prime, are exactly
\[
  1,\ 2,\ 4,\ 12.
\]
\end{theorem}

\begin{proof}
The numbers \(1,2,4,12\) are checked directly. For the converse, write
\[
  n = 2^k m
\]
with \(m\) odd. This is the standard odd-part decomposition and is the starting
point of the Lean proof.

If \(m=1\), then \(n\) is a power of two. The divisor \(8\) shows that
\(k\ge 3\) is impossible: indeed \(8\mid n\) would force \(9\mid n\) or \(9\)
prime, and neither can happen. Hence \(k\le 2\), so \(n\in\{1,2,4\}\).

Assume now \(m>1\). Taking \(d=m\), the number \(m+1\) is even and strictly
greater than \(2\), so it is not prime; therefore \(m+1\mid n\). Since
\(\gcd(m,m+1)=1\), this implies that \(m+1\) divides the power of two part,
hence
\[
  m+1 = 2^b
\]
for some \(b\le k\).

If \(b\) is odd, then \(2^b\mid n\). Applying the hypothesis to \(d=2^b\), we
would need \(2^b+1\mid n\) or \(2^b+1\) prime. But for odd \(b\), the number
\(2^b+1\) is divisible by \(3\), so it is composite, and it is also an odd
divisor larger than the odd part \(m=2^b-1\), so it cannot divide \(n\). This
is impossible. Thus \(b\) is even, say \(b=2c\).

If \(c=1\), then \(m=3\). The divisor \(3\) forces \(4\mid n\), and the same
``\(8\) implies \(9\)'' obstruction shows that \(k\) cannot exceed \(2\). Hence
\(n=12\).

It remains to rule out \(c\ge 2\). Since
\[
  m = 2^{2c}-1 = (2^c-1)(2^c+1),
\]
the divisor \(2^c+1\) also divides \(n\). Its successor \(2^c+2\) is even and
greater than \(2\), so it cannot be prime; therefore \(2^c+2\mid n\). Hence
the odd factor \(2^{c-1}+1\) divides \(n\), and therefore divides the odd part
\(m\). On the other hand,
\[
  m = (2^{c-1}+1)(2^{c+1}-4)+3,
\]
so \(2^{c-1}+1\mid 3\). This forces \(c=2\), hence \(m=15\). But then
\(24\mid n\), so the hypothesis at \(d=24\) would force \(25\mid n\) or \(25\)
prime. Since \(25\) is composite and is an odd divisor larger than the odd part
\(15\), this is impossible.

Therefore no further cases occur, and the only solutions are
\[
  n=1,2,4,12.
\]
\end{proof}
